% ===================================================================
%  Section I: Introduction
% ===================================================================

The deployment of reinforcement learning (RL) models in safety-critical
applications---autonomous driving, medical treatment planning, robotic
control---demands that stakeholders can verify the integrity of the training
process~\cite{levine2020offline}. Unlike supervised learning, where a fixed
mapping from inputs to labels can be checked through inference verification
alone, RL policies emerge from sequential interactions with an environment.
In the \emph{offline} setting, a policy is trained entirely from a fixed replay
dataset $\dataset$ without further environment interaction, making the
provenance and integrity of $\dataset$ as critical as the correctness of the
learning algorithm itself.

Consider a regulatory scenario where a company submits an offline-trained DQN
policy for certification. A dishonest party could alter a reward signal in the
replay buffer to make the policy appear better than it is, change a terminal
flag to hide catastrophic failure trajectories, substitute Bellman targets to
mask value-function errors, or tamper with gradient computations to conceal
unauthorized model modifications. Current proof-of-learning
frameworks~\cite{jia2021proofoflearning} address some of these concerns for
supervised training, but they do not capture the RL-specific semantics of replay
membership, Double-DQN target selection, temporal-difference (TD) error
computation, target-network synchronization, and checkpoint chaining that
characterize offline value-based methods.

Zero-knowledge proof systems offer a principled solution: a prover can
demonstrate that a computation was performed correctly without revealing
sensitive witness data~\cite{thalerbook}. Recent work has applied
zero-knowledge proofs to deep learning training~\cite{zkdl2023, kaizen2024,
zkpotvicinity2025} and inference~\cite{zkcnn2021, safetynets2017, vcnn2024},
but these systems target supervised architectures and do not address
RL-specific training semantics. At the same time, advances in general-purpose
zkVMs such as SP1~\cite{succinct2024sp1} and RISC Zero~\cite{risczero2023}
have made it feasible to prove complex Rust programs inside a STARK-based
virtual machine, opening the door to verifiable RL training.

This paper presents a \emph{layered artifact system} for zero-knowledge
verifiable offline DQN training. The system is deliberately scoped to selected
training fragments rather than claiming end-to-end training proofs, and it
separates concerns into distinct layers: dataset provenance, relation
semantics, proof generation, and proof aggregation. Each layer is independently
testable, and every claim is traceable to a concrete implementation artifact,
test, and benchmark row.

Our starting point is that RL does not merely add another model class to
this literature. A recent survey of zero-knowledge verifiable machine
learning~\cite{zkmlsurvey2026} organises proof-of-training work by
architecture---deep networks, convolutional networks, recurrent networks,
decision trees, language models---and lists no reinforcement-learning system
at all. The gap persists not because value-based RL is a larger computation,
but because it carries three verification obligations that supervised
proof-of-training constructions never encounter:

\begin{enumerate}[leftmargin=1.8em]
  \item \textbf{The labels are produced by the model being trained.} In
  supervised training the labels are fixed data, and committing to the dataset
  commits to them. A TD target is computed from the target network during
  training, so a proof must establish not only that the update used the label
  but that the label itself was derived correctly from a committed checkpoint.
  \item \textbf{Target-network synchronisation is a discrete event.} The
  target is refreshed every $N$ steps. A prover free to move those events
  changes what every subsequent label means, and the resulting run remains
  internally consistent.
  \item \textbf{Replay sampling has no canonical order.} Minibatch indices are
  drawn rather than enumerated, so a prover who selects the draw selects the
  data the claim is about.
\end{enumerate}

\noindent The third obligation is not hypothetical. During development our own
chain sampled with a fixed seed and a step index local to each fragment, so
every chunk of a $1248$-step run drew the same eight transitions from a
$50{,}000$-row dataset. Every hash agreed, every proof verified, and all $236$
tamper cases stayed green, because no field had been altered---the relation
simply described a different computation than intended. Binding the sampler to
the dataset commitment removes both this defect and the prover's freedom to
grind the seed~\cite{zkpotvicinity2025}.

The contributions of this work are as follows:
\begin{enumerate}[leftmargin=1.8em]
  \item \textbf{RL-specific verification relations.} We formalise and implement
  relations for the three obligations above---model-generated labels,
  synchronisation timing, and sampling integrity---together with replay
  membership and checkpoint chaining, and we bind the sampler to the committed
  dataset root so that minibatch selection follows from the commitment rather
  than from a prover input (\Cref{sec:relations}).

  \item \textbf{Proof-backed coverage with recursive aggregation.} We generate
  verified SP1 proofs for replay membership (to $100$k leaves), Bellman/TD
  arithmetic, forward evaluation, one-step SGD, checkpoint chains, training
  fragments, proof-manifest chunk chains for $T \in \{32,64,128\}$, and
  aggregation that verifies each child proof \emph{inside} the guest for
  $T \in \{16,32,64\}$ under a CUDA prover
  (\Cref{sec:proof-backend,sec:training-fragment}).

  \item \textbf{A whole training run under a single proof.} For both
  environments we prove a complete $1248$-step run as one root proof over a
  binary tree of eight $156$-step leaves, each leaf bound to the committed
  dataset the reported policy was trained on, with $\mathrm{step\_start}=0$,
  $\mathrm{step\_end}=1248$, and verification in $0.054$\,s
  (\Cref{sec:experiments}).

  \item \textbf{The provable configuration measured, not assumed.} The relation
  checks plain SGD on single transitions with component-wise gradient clipping,
  which is not what a tuned implementation runs. Rather than leave the gap to
  the reader we report it as its own row: on \texttt{lunarlander-random} the
  provable configuration reaches $-152.7$ against $-215.3$ for the tuned
  minibatch baseline, and an ablation isolates the target-sync interval---not
  the minibatch, and not the clipping rule the proof forces---as what governs
  whether it learns (\Cref{sec:experiments}).

  \item \textbf{Formal statements and an adversarial benchmark.} We state and
  sketch proofs for eight theorems spanning membership, dataset commitment,
  Bellman targets, update correctness, chain and aggregation soundness, and the
  zero-knowledge boundary, each mapped to relation code, SP1 backend, tests and
  benchmark rows; and we report tamper rejection across $19$ adversarial
  categories and $236$ cases with no unexpectedly accepted row
  (\Cref{sec:relations,sec:experiments}).
\end{enumerate}

\noindent\textbf{Scope boundaries.} This paper does not claim a complete proof
that an arbitrary DQN was trained from initialization to final model. It does
not claim Adam optimizer proofs, batch sizes above one, proof-backed
$k \in \{16,32,128\}$ fragments, or Atari/MuJoCo-scale proofs, and public
Minari/D4RL imports are source-integrity commitments rather than evidence of
honest collection. Proof length is bounded by cost rather than by arithmetic: a
CartPole run long enough to learn well would need roughly $68$ GPU-hours at the
leaf size used here. These boundaries are stated in \Cref{sec:discussion} and
in the formal theorem statements.

\smallskip
\noindent\textbf{Organization.} \Cref{sec:background} reviews background on
offline RL, DQN, and zero-knowledge proofs. \Cref{sec:threat-model} defines the
threat model and problem formulation. \Cref{sec:system-overview} describes the
system architecture. \Cref{sec:relations} presents the relations and formal
statements. \Cref{sec:proof-backend} details the SP1 proof backend.
\Cref{sec:training-fragment} covers training fragments and aggregation.
\Cref{sec:experiments} reports experimental results.
\Cref{sec:discussion} discusses scope and scalability.
\Cref{sec:related-work} surveys related work, and
\Cref{sec:conclusion} concludes with future directions.
